\section{Introduction}

\subsection{Complex Numbers}

Let's not start from math. Let's start from \textbf{physics / computation intuition}. Classical bits are represented by two distinct states: $0$ and $1$. In quantum computing, we have \textbf{qubits}, which are essentially \textbf{vector of numbers}. But not just any numbers: a qubit is a vector of \textbf{complex numbers}. They are needed because quantum systems have \textbf{both magnitude and phase}, and complex numbers can represent both. Real numbers can only represent ``\emph{how much}'' (magnitude), while complex numbers can represent ``\emph{how much}'' and ``\emph{in which direction / phase}'' (magnitude and phase).

\highspace
\textcolor{Green3}{\faIcon{question-circle} \textbf{Why complex (and not just real) numbers?}}
\begin{enumerate}
    \item In quantum computing, we don't store probabilities directly (e.g., $0.5$ for a 50\% chance). Instead, we store \textbf{probability amplitudes}, which are complex numbers. The probability of an outcome is the square of the magnitude of its amplitude:
    \begin{equation*}
        P = \left|z\right|^2
    \end{equation*}
    \item \textbf{Interference}. Quantum states can \textbf{combine and cancel}. For example, if we have two states with amplitudes $z_1$ and $z_2$, the combined amplitude is:
    \begin{equation*}
        z_{\text{combined}} = z_1 + z_2
    \end{equation*}
    If $z_1$ and $z_2$ have the same magnitude (\emph{how much}) but opposite phases (\emph{in which direction}), they can cancel each other out, leading to zero probability for that outcome. This is a key feature of quantum mechanics that allows for phenomena like quantum interference, which is essential for quantum algorithms.
\end{enumerate}
In general, complex numbers are needed because \textbf{quantum states behave like waves}, not like classical values. Waves have both amplitude and phase, and complex numbers are the natural way to represent them mathematically. This is why quantum mechanics and quantum computing rely heavily on complex numbers.

\highspace
\begin{flushleft}
    \textcolor{Green3}{\faIcon{question-circle} \textbf{What do we actually need?}}
\end{flushleft}
Not everything. Just a \textbf{minimal toolkit}. A complex number is:
\begin{equation*}
    z = x + iy
\end{equation*}
Where $x$ is the \textbf{real part} and $y$ is the \textbf{imaginary part}, and $i$ is the \textbf{imaginary unit} satisfying $i^2 = -1$. From a geometric interpretation, we can think of $z$ as a point in the 2D plane, where $x$ is the horizontal coordinate (real axis) and $y$ is the vertical coordinate (imaginary axis).

\newpage

\begin{flushleft}
    \textcolor{Green3}{\faIcon{book} \textbf{Polar form}}
\end{flushleft}
Instead of coordinates $(x, y)$, we can also represent a complex number in \textbf{polar form}:
\begin{equation*}
    z = re^{i \varphi}
\end{equation*}
Where:
\begin{itemize}
    \item $r = \left|z\right| = \sqrt{x^2 + y^2}$ is the \textbf{modulus} (or \emph{magnitude}) of the complex number, representing its \hl{distance from the origin in the complex plane.}
    \item $\varphi = \arg(z) = \tan^{-1} \left( \frac{y}{x} \right)$ is the \textbf{argument} (or \emph{phase} angle), representing the \hl{angle between the positive real axis and the line connecting the origin to the point $z$ in the complex plane.}
    \item The \textbf{exponential form} $e^{i \varphi}$ means that we can represent the complex number as a point on the unit circle in the complex plane, where $\varphi$ is the angle from the positive real axis. The modulus $r$ then scales this point to its actual distance from the origin.
    
    Using \hl{Euler's formula}:
    \begin{equation*}
        e^{i \varphi} = \cos \varphi + i \sin \varphi
    \end{equation*}
    We can rewrite $z$ as:
    \begin{equation*}
        z = r \left(\cos \varphi + i \sin \varphi\right)
    \end{equation*}
    This form is particularly useful in quantum computing because it makes it easy to understand how complex numbers can represent both magnitude and phase, which are essential for describing quantum states and their evolution.
\end{itemize}


\begin{figure}[!htp]
    \centering
    \begin{tikzpicture}[scale=1.2, >=Stealth]
        % Axes
        \draw[->, thick] (-0.5,0) -- (4.5,0) node[right] {real};
        \draw[->, thick] (0,-0.5) -- (0,3.5) node[above] {imaginary};

        % Point z
        \coordinate (O) at (0,0);
        \coordinate (Z) at (3,2);

        % Vector z
        \draw[thick, blue] (O) -- (Z) node[midway, above left] {$r=|z|$};

        % Projection on x-axis
        \draw[dashed] (Z) -- (3,0) node[below] {$x$};
        \draw[dashed] (Z) -- (0,2) node[left] {$y$};

        % Point label
        \filldraw[blue] (Z) circle (1.2pt) node[above right] {$z=x+iy=re^{i\varphi}$};

        % Angle phi
        \draw[->, thick, orange] (1,0) arc[start angle=0, end angle=33.69, radius=1];
        \node[orange] at (1.3,0.35) {$\varphi$};
    \end{tikzpicture}
    \caption{Geometric interpretation of a complex number in the complex plane. The arrow represents the complex number $z$ as a vector from the origin to the point $(x, y)$. The length of the vector is the modulus $r$, and the angle it makes with the positive real axis is the argument $\varphi$.}
\end{figure}

\newpage

\begin{flushleft}
    \textcolor{Green3}{\faIcon{book} \textbf{Complex Conjugate}}
\end{flushleft}
The \textbf{complex conjugate} of a complex number $z = x + iy$ is denoted as $\bar{z}$ and is defined as:
\begin{equation*}
    \bar{z} = x - iy
\end{equation*}
Geometrically, the complex conjugate $\bar{z}$ is the \textbf{reflection} of $z$ \textbf{across the real axis in the complex plane}. This means that if $z$ is represented by the point $(x, y)$, then $\bar{z}$ is represented by the point $(x, -y)$. The real part remains the same, while the imaginary part changes sign. See \autoref{fig:complex-conjugate} for a visual representation of this concept.

\begin{figure}[!htp]
    \centering
    \begin{tikzpicture}[scale=1.2, >=Stealth]

        % Axes
        \draw[->, thick] (-0.5,0) -- (4.5,0) node[right] {real};
        \draw[->, thick] (0,-3.5) -- (0,3.5) node[above] {imaginary};

        % Points
        \coordinate (O) at (0,0);
        \coordinate (Z) at (3,2);
        \coordinate (Zbar) at (3,-2);

        % Vectors
        \draw[thick, blue] (O) -- (Z) node[midway, above left] {$z$};
        \draw[thick, red] (O) -- (Zbar) node[midway, below left] {$\bar{z}$};

        % Points
        \filldraw[blue] (Z) circle (1.2pt) node[above right] {$x+iy$};
        \filldraw[red] (Zbar) circle (1.2pt) node[below right] {$x-iy$};

        % Projections
        \draw[dashed] (Z) -- (3,0);
        \draw[dashed] (Zbar) -- (3,0);

        % Vertical symmetry line
        \draw[dashed, gray] (Z) -- (Zbar);

        % Labels
        \node[below] at (3,0) {$x$};
        \node[left] at (0,2) {$y$};
        \node[left] at (0,-2) {$-y$};
    \end{tikzpicture}
    \caption{Geometric interpretation of the complex conjugate: reflection across the real axis.}
    \label{fig:complex-conjugate}
\end{figure}

\begin{flushleft}
    \textcolor{Green3}{\faIcon{book} \textbf{Magnitude}}
\end{flushleft}
The \textbf{magnitude} (or \textbf{modulus}) of a complex number $z = x + iy$ is denoted as $\left|z\right|$ and is defined as:
\begin{equation*}
    \left|z\right| = \sqrt{x^2 + y^2}
\end{equation*}
And it can also be expressed using the complex conjugate:
\begin{equation*}
    \left|z\right|^2 = z \cdot \bar{z} = (x + iy)(x - iy) = x^2 + y^2
\end{equation*}
The magnitude represents the \textbf{distance} of the point $z$ from the origin in the complex plane. It is \hl{always a non-negative real number.} The magnitude is important in quantum mechanics because the probability of an outcome is given by the square of the magnitude of the corresponding probability amplitude. For example, if a quantum state has an amplitude of $z$, the probability of measuring that state is $\left|z\right|^2$. This is why the magnitude of complex numbers is crucial in quantum computing, as it directly relates to the probabilities of different outcomes when measuring quantum states.

\highspace
\begin{flushleft}
    \textcolor{Green3}{\faIcon{sync} \textbf{Rotation in the complex plane}}
\end{flushleft}
Multiplying a complex number by a complex exponential of the form $e^{i\theta}$ corresponds to a \textbf{rotation} in the complex plane:
\begin{equation*}
    z' = z \cdot e^{i\theta} = r e^{i(\varphi + \theta)}
\end{equation*}
Geometrically, this operation:
\begin{itemize}
    \item \textbf{Preserves the magnitude} $|z|$
    \item \textbf{Adds $\theta$ to the phase} $\varphi$
\end{itemize}
Therefore, multiplication by $e^{i\theta}$ rotates the complex number by an angle $\theta$ around the origin.

\begin{figure}[!htp]
    \centering
    \begin{tikzpicture}[scale=1.2, >=Stealth]

        % Axes
        \draw[->, thick] (-0.5,0) -- (4.5,0) node[right] {real};
        \draw[->, thick] (0,-0.5) -- (0,3.5) node[above] {imaginary};

        % Origin
        \coordinate (O) at (0,0);

        % Original vector z (angle ~30°)
        \coordinate (Z) at ({3*cos(30)}, {3*sin(30)});

        % Rotated vector z' (angle ~70°)
        \coordinate (Zp) at ({3*cos(70)}, {3*sin(70)});

        % Draw vectors
        \draw[thick, blue] (O) -- (Z) node[midway, below right] {$z$};
        \draw[thick, red] (O) -- (Zp) node[midway, above left] {$z'$};

        % Points
        \filldraw[blue] (Z) circle (1.2pt);
        \filldraw[red] (Zp) circle (1.2pt);
        \node[blue] at (Z) [above right] {$re^{i\varphi}$};
        \node[red] at (Zp) [above] {$re^{i(\varphi+\theta)}$};

        % Arc for original angle phi
        \draw[->, thick, blue] (1,0) arc[start angle=0, end angle=30, radius=1];
        \node[blue] at (1.2,0.25) {$\varphi$};

        % Arc for rotation theta
        \draw[->, thick, orange] ({1.2*cos(30)}, {1.2*sin(30)})
            arc[start angle=30, end angle=70, radius=1.2];
        \node[orange] at (0.9,1.0) {$\theta$};

        % Radius label
        \node at (2,2) {$\left|z\right| = \left|z'\right|$};

    \end{tikzpicture}
    \caption{Rotation of a complex number: multiplying by $e^{i\theta}$ rotates the vector by angle $\theta$ while preserving its magnitude.}
\end{figure}

\highspace
\begin{flushleft}
    \textcolor{Green3}{\faIcon{exchange-alt} \textbf{Hermitian (Conjugate Transpose)}}
\end{flushleft}
Given a complex vector:
\begin{equation*}
    \mathbf{v} =
    \begin{bmatrix}
        a \\
        b
    \end{bmatrix}
\end{equation*}
Its \textbf{Hermitian} (or \textbf{conjugate transpose}) is defined as:
\begin{equation*}
    \mathbf{v}^\dagger =
    \begin{bmatrix}
        \bar{a} & \bar{b}
    \end{bmatrix}
\end{equation*}
This operation consists of:
\begin{itemize}
    \item Taking the \textbf{transpose} (column $\rightarrow$ row)
    \item Taking the \textbf{complex conjugate} of each element
\end{itemize}